\subsection{Bulk/boundary Modular Quintessence and DESI --- arXiv:2506.02731}

\textbf{Authors:} Luis A. Anchordoqui, Ignatios Antoniadis, Niccolo Cribiori, Arda Hasar, Dieter Lust, Joaquin Masias, Marco Scalisi

\paragraph{主な主張}
余剰次元半径 $R$ が暗黒エネルギーの大きさを、string coupling に対応する複素場 $S$ が時間発展を担う二場シナリオを提案し、species scale に基づく $SL(2,\mathbb{Z})$ 不変 potential による thawing quintessence が DESI DR2 と整合し得ることを示した \cite{Anchordoqui:2025modular}。

\paragraph{新規性}
self-dual quintessence の発想を full modular invariance へ拡張し、$-\log[|\eta(S)|^4(S+\bar S)]$ で表される modular-invariant species count を通じて potential の形に quantum-gravity 的な動機を与えた点にある。

\paragraph{位置付け}
物理的に動機づけられた modular-invariant potential の具体例として有用だが、有限の modular vacuum への capture や stabilization 自体は解いていない。そのため dynamical stabilization の直接の先行研究というより、potential 選択の指針・benchmark として重要である。
